% Chapter: Mathlib Integration — The Translation Functor
% Part of the CCTT Monograph — Proof Calculus, Chapter IV

\chapter{Mathlib Integration --- The Translation Functor}
\label{ch:c4-mathlib}
\label{ch:mathlib-integration}

A key design principle of $\mathrm{C}^4$ is that any theorem proved in
Lean~4's Mathlib library can be imported as a trusted axiom.  This chapter
defines the translation functor $\Phi : \mathrm{CIC}_{\mathrm{Lean}} \to
\mathrm{C}^4$, gives the concrete algorithm for translating theorem
\emph{statements} (not proofs), proves soundness, and works through
examples from the 177,744-theorem Mathlib catalog.


% ─────────────────────────────────────────────────────────
\section{Design: Statements, Not Proofs}
\label{sec:statements-not-proofs}

The translation functor imports \emph{theorem types}, not proof terms.
This is a deliberate design choice:

\begin{enumerate}
\item \textbf{Lean proofs are large.}  A typical Mathlib proof term is
  thousands of nodes.  Translating the full proof would be expensive and
  fragile.

\item \textbf{Lean is trusted.}  Lean's kernel has been extensively
  verified.  If Lean accepts a proof of $P$, we trust that $P$ is
  true (relative to Lean's axioms: classical logic + quotient types +
  propositional extensionality + function extensionality + choice).

\item \textbf{$\mathrm{C}^4$ is conservative over CIC.}
  By Theorem~\ref{thm:conservative-extension}, any CIC-provable
  statement remains provable in $\mathrm{C}^4$.  Therefore, importing
  Lean's theorem statements as axioms does not introduce inconsistency
  (assuming Lean is consistent).

\item \textbf{Economy.}  We need the theorem's \emph{interface}
  (its type), not its implementation (its proof term), just as in
  programming we use a library's API without reading its source code.
\end{enumerate}

The translated statements enter $\mathrm{C}^4$ with trust level
$\mathsf{LEAN\_VERIFIED}$---the highest level in the CCTT trust
hierarchy.


% ─────────────────────────────────────────────────────────
\section{The Translation Functor $\Phi$}
\label{sec:translation-functor}

\begin{definition}[Translation Functor]
\label{def:translation-functor}
The \textbf{translation functor} $\Phi : \mathrm{CIC}_{\mathrm{Lean}}
\to \mathrm{C}^4$ maps:
\begin{enumerate}
\item \textbf{Universes:} $\mathsf{Sort}\, u \mapsto \mathcal{U}_u$
  (with $\mathsf{Prop} \mapsto \mathsf{Prop}$).
\item \textbf{Types:} Lean types to $\mathrm{C}^4$ types (see the
  type mapping table below).
\item \textbf{Terms:} Lean terms (in types) to $\mathrm{C}^4$ terms.
\item \textbf{Declarations:}
  $\mathsf{theorem}\; n : P := t \;\mapsto\;
  \mathsf{axiom}\; \Phi(n) : \Phi(P)$.
\end{enumerate}
$\Phi$ is a functor in the categorical sense: it preserves identity
and composition of morphisms in the type-theoretic category.
\end{definition}


% ─────────────────────────────────────────────────────────
\section{The Translation Algorithm}
\label{sec:translation-algorithm}

The algorithm is implemented in the code modules
\texttt{lean\_parser.py}, \texttt{lean\_type\_mapper.py}, and
\texttt{prop\_translator.py} in \texttt{new\_src/cctt/mathlib\_bridge/}.
Here we give the formal specification.

\subsection{Step 1: Parse the Lean~4 Declaration}

Given a Lean~4 declaration string, the parser (a recursive-descent
parser for Lean~4 syntax) produces an AST node.  The relevant fields are:
\begin{itemize}
\item \textbf{Name}: the fully qualified theorem name (e.g.,
  \texttt{Nat.add\_comm}).
\item \textbf{Universe parameters}: $u_1, \ldots, u_k$.
\item \textbf{Type}: the proposition type $P$ (a Lean expression).
\item \textbf{Proof term}: discarded.
\end{itemize}

\subsection{Step 2: Extract the Proposition Type}

From the AST, extract $P$ and discard the proof term.  If the
declaration is a \texttt{def} rather than a \texttt{theorem}, extract
the return type.

\subsection{Step 3: Map Lean Types to $\mathrm{C}^4$ Types}

The type mapping $\Phi_\mathrm{type}$ is defined by structural recursion
on the Lean type:

\begin{center}
\begin{tabular}{lll}
\toprule
\textbf{Lean Type} & \textbf{$\mathrm{C}^4$ Type} & \textbf{Notes} \\
\midrule
$\mathsf{Nat}$ & $\Nat$ & natural numbers \\
$\mathsf{Int}$ & $\Int$ & integers \\
$\mathsf{Bool}$ & $\Bool$ & booleans \\
$\mathsf{String}$ & $\mathsf{String}$ & strings \\
$\mathsf{List}\;\alpha$ & $\mathsf{List}(\Phi(\alpha))$ & polymorphic \\
$\mathsf{Array}\;\alpha$ & $\mathsf{List}(\Phi(\alpha))$ & arrays $\to$ lists \\
$\mathsf{Option}\;\alpha$ & $\mathsf{Option}(\Phi(\alpha))$ & \\
$\mathsf{Fin}\;n$ & $\{k : \Nat \mid k < n\}$ & bounded nat \\
$\mathsf{Prod}\;\alpha\;\beta$ & $\Phi(\alpha) \times \Phi(\beta)$ & \\
$\mathsf{Sum}\;\alpha\;\beta$ & $\Phi(\alpha) + \Phi(\beta)$ & \\
$\mathsf{Prop}$ & $\mathsf{Prop}$ & propositions \\
$\mathsf{Sort}\;u$ & $\mathcal{U}_u$ & universes \\
\bottomrule
\end{tabular}
\end{center}

For composite types, the mapping recurses:
\begin{itemize}
\item $\Phi(\Pi(x:\alpha).\,\beta) = \Pi(x:\Phi(\alpha)).\,\Phi(\beta)$
\item $\Phi(\Sigma(x:\alpha).\,\beta) = \Sigma(x:\Phi(\alpha)).\,\Phi(\beta)$
\item $\Phi(\alpha \to \beta) = \Phi(\alpha) \to \Phi(\beta)$
  \quad (non-dependent $\Pi$)
\end{itemize}

\subsection{Step 4: Map Lean Type Classes to $\mathrm{C}^4$ Assumptions}

Lean type classes become explicit assumptions in $\mathrm{C}^4$:

\begin{center}
\begin{tabular}{ll}
\toprule
\textbf{Lean Type Class} & \textbf{$\mathrm{C}^4$ Assumption} \\
\midrule
$[\mathsf{Add}\;\alpha]$ & $\alpha$ has operation $+ : \alpha \to \alpha \to \alpha$ \\
$[\mathsf{Mul}\;\alpha]$ & $\alpha$ has operation $\times : \alpha \to \alpha \to \alpha$ \\
$[\mathsf{Ord}\;\alpha]$ & $\alpha$ has a linear order \\
$[\mathsf{CommMonoid}\;\alpha]$ & $(\alpha, \cdot, 1)$ is a commutative monoid \\
$[\mathsf{DecidableEq}\;\alpha]$ & equality on $\alpha$ is decidable \\
$[\mathsf{Fintype}\;\alpha]$ & $\alpha$ is finite \\
\bottomrule
\end{tabular}
\end{center}

\noindent
Formally: $\Phi([\mathsf{C}\;\alpha]\;\beta) =
\Pi(\mathsf{inst} : \Phi(\mathsf{C})(\Phi(\alpha))).\,\Phi(\beta)$.

\subsection{Step 5: Translate Quantifiers and Connectives}

\begin{center}
\begin{tabular}{ll}
\toprule
\textbf{Lean Connective} & \textbf{$\mathrm{C}^4$ Translation} \\
\midrule
$\forall (x : \alpha), P$ & $\Pi(x : \Phi(\alpha)).\,\Phi(P)$ \\
$\exists (x : \alpha), P$ & $\Sigma(x : \Phi(\alpha)).\,\Phi(P)$ \\
$P \wedge Q$ & $\Phi(P) \times \Phi(Q)$ \\
$P \vee Q$ & $\Phi(P) + \Phi(Q)$ \\
$P \to Q$ & $\Phi(P) \to \Phi(Q)$ \\
$\neg P$ & $\Phi(P) \to \bot$ \\
$P \leftrightarrow Q$ & $(\Phi(P) \to \Phi(Q)) \times (\Phi(Q) \to \Phi(P))$ \\
$a = b$ & $\Path_{\Phi(A)}\,\Phi(a)\,\Phi(b)$ \\
\bottomrule
\end{tabular}
\end{center}

\noindent
The translation of equality to path types is the key design choice:
Lean's propositional equality becomes a cubical path in $\mathrm{C}^4$.

\subsection{Step 6: Handle Universe Polymorphism}

Lean theorems are often polymorphic in universe levels.  We translate
$\mathsf{Sort}\;u$ to $\mathcal{U}_u$ and instantiate universe variables
to concrete levels when needed.  For most Mathlib theorems about $\Nat$,
$\Int$, $\mathsf{List}$, etc., the universe level is $0$ or $1$.

\subsection{Step 7: Emit as $\mathrm{C}^4$ Axiom}

The final output is:
\[
  \mathsf{axiom}\;\Phi(\mathit{name}) : \Phi(\mathit{type})
  \quad [\mathsf{trust} = \mathsf{LEAN\_VERIFIED}]
\]


% ─────────────────────────────────────────────────────────
\section{Soundness of the Translation}
\label{sec:translation-soundness}

\begin{theorem}[Soundness of $\Phi$]
\label{thm:phi-soundness}
If Lean proves $\vdash_{\mathrm{Lean}} t : P$, then the axiom
$\Phi(P)$ is consistent in $\mathrm{C}^4$ (i.e., adding $\Phi(P)$
as an axiom does not make $\mathrm{C}^4$ inconsistent, assuming Lean
is consistent).
\end{theorem}

\begin{proof}
We argue in three steps.

\textbf{Step 1: $\Phi$ preserves types.}
By structural induction on Lean types, we show that $\Phi$ maps
well-formed Lean types to well-formed $\mathrm{C}^4$ types.  The key
cases:
\begin{itemize}
\item $\Phi(\Nat) = \Nat$ is a type in $\mathcal{U}_0$.
\item $\Phi(\Pi(x:\alpha).\beta) = \Pi(x:\Phi(\alpha)).\Phi(\beta)$
  is well-formed by the $\Pi$-formation rule, given that
  $\Phi(\alpha)$ and $\Phi(\beta)$ are well-formed (IH).
\item $\Phi(a =_A b) = \Path_{\Phi(A)}\,\Phi(a)\,\Phi(b)$ is
  well-formed by the $\Path$-formation rule.
\item Type classes: $\Phi$ turns implicit class arguments into
  explicit $\Pi$-bound assumptions, which is type-correct.
\end{itemize}

\textbf{Step 2: Conservativity.}
By Theorem~\ref{thm:conservative-extension}, $\mathrm{C}^4$ is
conservative over CIC.  The CIC fragment of $\mathrm{C}^4$ includes
all of Lean's core type theory (CIC with universes, $\Pi$, $\Sigma$,
inductive types).  Therefore, if a proposition $P$ is provable in
Lean (CIC + axioms), then $\Phi(P)$ is provable in the CIC fragment
of $\mathrm{C}^4$ (since $\Phi$ acts as the identity on CIC-level
constructs up to the equality translation).

\textbf{Step 3: Equality translation.}
The translation $a =_A b \mapsto \Path_{\Phi(A)}\,\Phi(a)\,\Phi(b)$
uses path types instead of inductive identity.  By univalence
(Theorem~\ref{thm:univalence}) and the fact that identity types and
path types are equivalent in cubical type theory, this preserves
provability: every proof of $a = b$ in Lean can be turned into a
path $\langle i \rangle\, t : \Path\,a\,b$ in $\mathrm{C}^4$.

Therefore, $\Phi(P)$ is inhabited in $\mathrm{C}^4$ whenever $P$
is proved in Lean, so adding $\Phi(P)$ as an axiom cannot introduce
inconsistency (it is already provable).
\end{proof}


% ─────────────────────────────────────────────────────────
\section{Concrete Examples}
\label{sec:concrete-examples}

We now translate several real Mathlib theorems.

\begin{example}[\texttt{Nat.add\_comm}]
\textbf{Lean declaration:}
\begin{lstlisting}[language={}]
theorem Nat.add_comm (n m : Nat) : n + m = m + n
\end{lstlisting}
\textbf{$\mathrm{C}^4$ axiom:}
\[
  \mathsf{axiom}\;\Phi(\mathsf{Nat.add\_comm}) :
  \Pi(n\, m : \Nat).\, \Path_\Nat\, (n + m)\, (m + n)
\]
\end{example}

\begin{example}[\texttt{List.map\_map}]
\textbf{Lean declaration:}
\begin{lstlisting}[language={}]
theorem List.map_map (g : B -> C) (f : A -> B) (l : List A) :
  List.map g (List.map f l) = List.map (g . f) l
\end{lstlisting}
\textbf{$\mathrm{C}^4$ axiom:}
\[
  \mathsf{axiom}\;\Phi(\mathsf{List.map\_map}) :
  \Pi(A\, B\, C : \mathcal{U}_0).\,
  \Pi(g : B \to C).\, \Pi(f : A \to B).\,
  \Pi(l : \mathsf{List}(A)).\,
  \Path_{\mathsf{List}(C)}\,
    (\mathsf{map}\;g\;(\mathsf{map}\;f\;l))\,
    (\mathsf{map}\;(g \circ f)\;l)
\]
\end{example}

\begin{example}[\texttt{Finset.sum\_comm}]
\textbf{Lean declaration:}
\begin{lstlisting}[language={}]
theorem Finset.sum_comm [AddCommMonoid M]
  (s : Finset A) (t : Finset B) (f : A -> B -> M) :
  s.sum (fun a => t.sum (f a)) = t.sum (fun b => s.sum (fun a => f a b))
\end{lstlisting}
\textbf{$\mathrm{C}^4$ axiom:}
\begin{align*}
  &\mathsf{axiom}\;\Phi(\mathsf{Finset.sum\_comm}) : \\
  &\quad \Pi(A\, B\, M : \mathcal{U}_0).\,
  \Pi(\mathsf{inst} : \mathsf{AddCommMonoid}(M)).\, \\
  &\quad \Pi(s : \mathsf{Finset}(A)).\,
  \Pi(t : \mathsf{Finset}(B)).\,
  \Pi(f : A \to B \to M).\, \\
  &\quad \Path_M\;
    (\mathsf{sum}(s, \lambda a.\, \mathsf{sum}(t, f\,a)))\;
    (\mathsf{sum}(t, \lambda b.\, \mathsf{sum}(s, \lambda a.\, f\,a\,b)))
\end{align*}
\end{example}

\begin{example}[\texttt{Nat.sub\_add\_cancel}]
\textbf{Lean:}
\begin{lstlisting}[language={}]
theorem Nat.sub_add_cancel {n m : Nat} (h : m <= n) : n - m + m = n
\end{lstlisting}
\textbf{$\mathrm{C}^4$:}
\[
  \mathsf{axiom}\;\Phi(\mathsf{Nat.sub\_add\_cancel}) :
  \Pi(n\, m : \Nat).\, (m \leq n) \to
  \Path_\Nat\, ((n - m) + m)\, n
\]
\end{example}

\begin{example}[\texttt{List.length\_map}]
\textbf{Lean:}
\begin{lstlisting}[language={}]
theorem List.length_map (f : A -> B) (l : List A) :
  (List.map f l).length = l.length
\end{lstlisting}
\textbf{$\mathrm{C}^4$:}
\[
  \mathsf{axiom}\;\Phi(\mathsf{List.length\_map}) :
  \Pi(A\, B : \mathcal{U}_0).\,
  \Pi(f : A \to B).\, \Pi(l : \mathsf{List}(A)).\,
  \Path_\Nat\, (\mathsf{length}(\mathsf{map}\;f\;l))\,
    (\mathsf{length}(l))
\]
\end{example}

\begin{example}[\texttt{Nat.mul\_comm}]
\textbf{Lean:}
\begin{lstlisting}[language={}]
theorem Nat.mul_comm (n m : Nat) : n * m = m * n
\end{lstlisting}
\textbf{$\mathrm{C}^4$:}
\[
  \mathsf{axiom}\;\Phi(\mathsf{Nat.mul\_comm}) :
  \Pi(n\, m : \Nat).\, \Path_\Nat\, (n \times m)\, (m \times n)
\]
\end{example}

\begin{example}[\texttt{List.reverse\_reverse}]
\textbf{Lean:}
\begin{lstlisting}[language={}]
theorem List.reverse_reverse (l : List A) :
  l.reverse.reverse = l
\end{lstlisting}
\textbf{$\mathrm{C}^4$:}
\[
  \mathsf{axiom}\;\Phi(\mathsf{List.reverse\_reverse}) :
  \Pi(A : \mathcal{U}_0).\, \Pi(l : \mathsf{List}(A)).\,
  \Path_{\mathsf{List}(A)}\,
    (\mathsf{reverse}(\mathsf{reverse}(l)))\, l
\]
\end{example}


% ─────────────────────────────────────────────────────────
\section{OTerm Mapping}
\label{sec:oterm-mapping}

When a Mathlib theorem concerns $\Nat$, $\Int$, $\mathsf{List}$, or other
types that have OTerm counterparts, we can additionally map the theorem
to an OTerm equality in $\mathcal{U}_\PyObj$.

\begin{definition}[OTerm-Compatible Theorem]
A theorem $\Phi(P)$ is \textbf{OTerm-compatible} if every type in $P$
maps to an OTerm type fiber:
\begin{center}
\begin{tabular}{ll}
$\Nat$ & $\to \mathsf{TInt}$ fiber (non-negative) \\
$\Int$ & $\to \mathsf{TInt}$ fiber \\
$\mathsf{List}(\alpha)$ & $\to \mathsf{TList}$ fiber \\
$\mathsf{String}$ & $\to \mathsf{TStr}$ fiber \\
$\mathsf{Bool}$ & $\to \mathsf{TBool}$ fiber \\
\end{tabular}
\end{center}
\end{definition}

For OTerm-compatible theorems, $\Phi(P)$ can be used as both a
$\mathrm{C}^4$ axiom and a rewrite rule in the OTerm universe, bridging
the gap between Lean's mathematical library and CCTT's program reasoning.


% ─────────────────────────────────────────────────────────
\section{The Theorem Catalog}
\label{sec:catalog}

The file \texttt{mathlib\_full\_catalog.json} contains 177,744 theorem
entries extracted from Mathlib4.  Each entry has:

\begin{itemize}
\item \textbf{name}: the fully qualified Lean name.
\item \textbf{type}: the Lean type expression (as a string).
\item \textbf{module}: the Mathlib module path.
\item \textbf{docstring}: documentation (if any).
\end{itemize}

\subsection{Querying the Catalog}

To use a Mathlib theorem in a $\mathrm{C}^4$ proof:

\begin{enumerate}
\item \textbf{Search by name:} Look up the theorem by its Lean name
  (e.g., \texttt{Nat.add\_comm}).  The catalog is indexed by name for
  $O(1)$ lookup.

\item \textbf{Search by type pattern:} For ``find me a theorem about
  list length after map,'' search for type strings containing
  \texttt{List.map} and \texttt{length}.

\item \textbf{Translate:} Apply the functor $\Phi$ to the theorem type.

\item \textbf{Instantiate:} Substitute concrete types and values for
  universe and type variables.

\item \textbf{Use in proof:} Apply the resulting $\mathrm{C}^4$ axiom
  via the $\mathsf{MathlibTheorem}$ proof term constructor.
\end{enumerate}

\subsection{Catalog Statistics}

\begin{center}
\begin{tabular}{lr}
\toprule
\textbf{Category} & \textbf{Count} \\
\midrule
Total theorems & 177,744 \\
About $\Nat$ operations & $\sim$4,200 \\
About $\mathsf{List}$ operations & $\sim$3,100 \\
About $\Int$ operations & $\sim$2,800 \\
About algebraic structures & $\sim$18,500 \\
About topology & $\sim$12,300 \\
About analysis & $\sim$15,600 \\
OTerm-compatible (estimated) & $\sim$12,000 \\
\bottomrule
\end{tabular}
\end{center}


% ─────────────────────────────────────────────────────────
\section{Integration with the Proof System}
\label{sec:mathlib-proof-integration}

\begin{definition}[MathlibTheorem Proof Term]
The proof term $\mathsf{MathlibTheorem}(\mathit{name}, \sigma)$
applies a translated Mathlib theorem with instantiation $\sigma$:
\[
  \frac{\mathit{name} \in \mathrm{Catalog} \qquad
        \sigma : \mathrm{FV}(\Phi(\mathit{type})) \to \mathrm{Terms}}
       {\Gamma \vdash \mathsf{MathlibTheorem}(\mathit{name}, \sigma) :
        \Phi(\mathit{type})[\sigma]}
\]
\end{definition}

\begin{remark}
The $\mathsf{MathlibTheorem}$ constructor is already present in the
proof term language (Chapter~\ref{ch:proof-theory-foundations}).  The
contribution of this chapter is the \emph{translation functor} that
makes the theorem type inhabitable in $\mathrm{C}^4$, together with
the soundness guarantee.
\end{remark}

\begin{example}[Using Mathlib in an Equivalence Proof]
To prove that $\mathsf{sum}(l_1 \mathbin{+\!\!+} l_2) =
\mathsf{sum}(l_1) + \mathsf{sum}(l_2)$ for lists, we can use:
\begin{enumerate}
\item $\Phi(\mathsf{List.sum\_append})$ from Mathlib, which gives
  $\Path\, (\mathsf{sum}(l_1 \mathbin{+\!\!+} l_2))\,
  (\mathsf{sum}(l_1) + \mathsf{sum}(l_2))$ directly.
\item Alternatively, prove it by list induction + axiom D1 (commutativity
  of addition).  But the Mathlib import is shorter and more trustworthy.
\end{enumerate}
\end{example}
